\subsection{Az alkalmazás áttekintése}

A fejlesztett alkalmazás elsősorban Android operációs rendszerre készült. 
A projekt alapját a Kotlin Multiplatform adja, amely lehetővé teszi, hogy az alkalmazás üzleti 
logikájának és adatkezelésének egy része platformfüggetlen módon készüljön el \cite{kotlinmultiplatform}.
Ennek előnye, hogy a közös kódrészek a későbbiekben más platformokon, például Androidon és iOS-en is
felhasználhatók. Az iOS-verzió tényleges implementálása és tesztelése azonban teszteszköz hiányában
ezen projekt során nem történt meg, ezért a dolgozat elsősorban az Androidos megvalósításra fókuszál.

A fejlesztés nyelve Kotlin volt, amely az Android-fejlesztésben széles körben támogatott, modern 
programozási nyelv. A Java nyelvvel való együttműködése miatt jól illeszkedik az Android környezetbe,
ugyanakkor számos olyan nyelvi funkciót kínál, amely megkönnyíti a biztonságosabb és átláthatóbb kód 
írását, például a tömörebb szintaxis és az adatosztályok használata. \cite{kotlinandroid}

A felhasználói felület (UI) a Jetpack Compose keretrendszer komponenseiből épül fel. A hagyományos,
imperatív (XML-alapú) Androidos nézetekkel szemben a Compose deklaratív megközelítést alkalmaz: az egyes
képernyők (Composable függvények) automatikusan és dinamikusan újrarajzolják magukat (recomposition),
amint az alkalmazás mögöttes állapota megváltozik \cite{compose}.


\subsubsection{A rendszer felépítése és folyamata}

A rendszer felépítése több, egymástól elkülönített komponensből áll :
\begin{itemize}
    \item A felhasználói felületet megvalósító képernyők 
    \item A munkamenet állapotát kezelő logika 
    \item A szenzoradatok fogadásáért és feldolgozásáért felelős részek 
    \item Adattárolási és visszakeresési mechanizmusok 
    \item Eredmények megjelenítésére szolgáló komponensek 
\end{itemize}

Ezek a komponensek együttműködve alkotják meg a teljes alkalmazást, és a következő felhasználói
folyamatot valósítják meg :
\begin{enumerate}
    \item Az alkalmazás indítása, mely során a főképernyő jelenik meg.
    \item A Sensoria eszköz csatlakoztatása Bluetooth Low Energy (BLE) protokollon keresztül,
    és a nyers szenzoradatok fogadása.
    \item A gyakorlatok listájának megjelenítése, amely lehet előre meghatározott vagy a
    felhasználó által kiválasztott gyakorlatokat tartalmazhat.
    \item A mérés indítása.
    \item A szenzoradatok valós idejű fogadása, illetve a gyakorlatok végrehajtásának videós megjelenítése.
    \item A mérés lezárása, amely során a gyakorlatokhoz kapcsolódó adatok mentése történik.
    \item A mérés eredményeinek megjelenítése, amely lehetőséget ad a gyakorlatok áttekintésére és összehasonlítására.
    \item A korábban rögzített munkamenetek visszakeresése és megtekintése.
\end{enumerate}

A dokumentáció a további fejezetekben ezeknek a komponenseknek a részletes ismertetésére,
valamint a fejlesztés során felmerülő kihívásokra és megoldásokra fókuszál.